\chapter*{Abstract}
\addcontentsline{toc}{chapter}{Abstract}
\section*{Rezumat}

Lucrarea de față propune un sistem de analiză și predicție a probabilității
de câștig în meciuri de fotbal, bazat pe tehnici de machine learning. Sistemul
procesează evenimentele unui meci după încheierea acestuia, reconstituind
evoluția probabilităților de-a lungul partidei pe baza evenimentelor
înregistrate: goluri, șuturi, expected goals (xG), presiune și cartonașe.
Arhitectura adoptată este proiectată astfel încât să permită extinderea
naturală către predicție în timp real, fără modificări structurale majore.

Un model \textbf{LightGBM} a fost antrenat pe \textbf{12.2 milioane de
evenimente} provenite din \textbf{3.464 de meciuri} din setul de date public
StatsBomb, acoperind competițiile majore europene din perioada 2017--2023.
Modelul obține un \textbf{RPS de 0.1431}, un \textbf{Brier Score de 0.1472}
și o \textbf{eroare de calibrare de 0.0309}, valori competitive față de
literatura de specialitate și sub pragul academic de 0.05 pentru uz tactic
real \cite{Robberechts2019}.

Pe lângă predicție, sistemul identifică automat momentele decisive ale
meciului pe baza variațiilor semnificative de probabilitate, oferind
analiștilor un instrument concret pentru înțelegerea dinamicii unui joc.
Aplicația este implementată ca un sistem full-stack: un backend REST API
construit cu FastAPI, autentificare securizată prin JWT și un frontend React,
care oferă acces la clasamentul și meciurile Ligii 1 România (Superliga),
sezonul 2023, prin integrarea API-Football. Utilizatorii pot genera predicții
probabilistice pentru meciurile finalizate, predicțiile fiind salvate în
istoricul personal al fiecărui utilizator.

\bigskip
\noindent\textbf{Cuvinte cheie:} machine learning, win probability, fotbal,
LightGBM, predicție post-meci, StatsBomb, expected goals, calibrare
probabilistică.

% ----------------------------------------------------------------
\section*{Abstract (English)}

This thesis presents a system for analysing and predicting win probability
in football matches using machine learning. The system processes match events
after the game has concluded, reconstructing the evolution of win probabilities
throughout the match based on recorded events: goals, shots, expected goals
(xG), pressure, and cards. The architecture is designed to allow natural
extension towards live, in-game prediction without major structural changes.

A \textbf{LightGBM} model was trained on \textbf{12.2 million events} from
\textbf{3,464 matches} using the StatsBomb Open Data dataset, covering major
European competitions from 2017 to 2023. The model achieves an
\textbf{RPS of 0.1431}, a \textbf{Brier Score of 0.1472}, and a
\textbf{Calibration Error of 0.0309}, results competitive with the
existing literature and below the 0.05 threshold considered reliable for
tactical use \cite{Robberechts2019}.

Beyond prediction, the system automatically identifies key moments in a
match by detecting significant probability shifts, providing analysts with
a concrete tool for understanding match dynamics.

The application is implemented as a full-stack system comprising a FastAPI
REST backend with JWT authentication and a React frontend. It integrates
the API-Football service to provide access to the Romanian Liga 1 (Superliga)
standings and fixtures from the 2023 season. Users can generate probabilistic
predictions for completed matches, with each prediction saved to their
personal history.

\bigskip
\noindent\textbf{Keywords:} machine learning, win probability, football,
LightGBM, post-match analysis, StatsBomb, expected goals,
probabilistic calibration.

% ----------------------------------------------------------------
\chapter*{Introducere}
\addcontentsline{toc}{chapter}{Introducere}
\label{intro}

\section*{Prezentarea temei}

Fotbalul este unul dintre cele mai analizate sporturi din lume, iar
interesul pentru cuantificarea performanței și pentru predicția
rezultatelor a crescut semnificativ odată cu disponibilitatea datelor
granulare la nivel de eveniment. Totuși, majoritatea soluțiilor
existente se concentrează fie pe predicția pre-meci, fie pe sisteme
comerciale care nu sunt accesibile publicului larg sau mediului academic.

Această lucrare abordează problema estimării probabilității de câștig
\textit{în timpul} unui meci de fotbal, pe baza evenimentelor înregistrate
(goluri, șuturi, expected goals (xG), presiune și cartonașe),
reconstituind evoluția probabilistică a partidei eveniment cu eveniment.
Spre deosebire de predicția clasică pre-meci, această abordare surprinde
dinamica reală a jocului și permite identificarea momentelor decisive
care au influențat decisiv soarta meciului.

Relevanța practică a temei este dublă: pe de o parte, analiștii sportivi
și staff-urile tehnice pot utiliza astfel de instrumente pentru evaluarea
tactică post-meci; pe de altă parte, arhitectura propusă este proiectată
explicit să permită extensia naturală către predicție în timp real, fără
modificări structurale majore.

\section*{Obiectivele lucrării}

Principalele obiective urmărite în cadrul acestei lucrări sunt:

\begin{enumerate}
    \item Construirea unui pipeline de procesare a datelor la nivel de
    eveniment, pornind de la setul de date public StatsBomb Open Data,
    acoperind \textbf{3.464 de meciuri} și \textbf{12.2 milioane de
    evenimente} din competițiile majore europene (2017--2023).

    \item Antrenarea și evaluarea unui model \textbf{LightGBM} capabil
    să estimeze probabilitățile de câștig, egal și pierdere la fiecare
    moment al meciului, cu metrici competitive față de literatura de
    specialitate.

    \item Implementarea unui mecanism automat de detectare a momentelor
    decisive ale meciului, bazat pe variațiile semnificative ale
    probabilităților estimate.

    \item Dezvoltarea unei aplicații full-stack funcționale care expune
    modelul printr-un REST API (FastAPI, autentificare JWT) și un frontend
    React, integrat cu API-Football pentru accesul la datele Ligii 1
    România (Superliga), sezonul 2023.

    \item Validarea rezultatelor prin metrici probabilistice standard:
    Ranked Probability Score (RPS), Brier Score și eroare de calibrare,
    cu raportare la pragurile acceptate în literatura de specialitate
    \cite{Robberechts2019}.
\end{enumerate}

\section*{Contribuția proprie și rezultate originale}

Contribuția principală a acestei lucrări constă în proiectarea și
implementarea integrală a unui sistem end-to-end de predicție a
probabilității de câștig în fotbal, de la procesarea brută a datelor
până la interfața utilizatorului final.

Din punct de vedere tehnic, contribuțiile originale includ:

\begin{itemize}
    \item Construirea unui set de caracteristici temporale derivate din
    evenimentele StatsBomb, care codifică starea meciului la fiecare
    moment: scor curent, presiune cumulată, xG acumulat pe echipă,
    avantaj numeric și timp rămas.

    \item Antrenarea unui model LightGBM pe acest set de date, obținând
    un \textbf{RPS de 0.1431}, un \textbf{Brier Score de 0.1472} și o
    \textbf{eroare de calibrare de 0.0309}, valori sub pragul de 0.05
    considerat acceptabil pentru aplicații tactice reale
    \cite{Robberechts2019}. Eroarea de calibrare obținută se situează
    semnificativ sub acest prag, indicând că probabilitățile produse de
    model reflectă fidel frecvențele reale.

    \item Un algoritm de detectare automată a momentelor decisive, care
    marchează punctele din meci în care probabilitatea de câștig a
    înregistrat cele mai abrupte variații, oferind o naratologie
    cantitativă a partidei.

    \item O arhitectură software modulară, cu separare clară între
    componenta de ML, backend-ul API și frontend, astfel încât
    înlocuirea modelului sau extinderea la predicție live să nu
    implice rescrierea sistemului.
\end{itemize}

\section*{Structura lucrării}

Lucrarea este organizată în următoarele capitole:

\textbf{Capitolul 1 --- Fundamente teoretice și studiu de literatură}
prezintă modelele de învățare automată utilizate (Gradient Boosting,
LightGBM, Regresie Logistică) și o sinteză a literaturii relevante privind
modelarea win probability în sport.

\textbf{Capitolul 2 --- Metodologie și experimente} descrie formularea
problemei, setul de date StatsBomb Open Data, construirea caracteristicilor
temporale, strategia de experimentare și optimizarea hiperparametrilor,
metricile de evaluare probabilistică (RPS, Brier Score, eroare de calibrare),
comparația celor trei modele și selecția modelului final LightGBM, precum și
algoritmul de detectare a momentelor decisive.

\textbf{Capitolul 3 --- Aplicația dezvoltată} descrie arhitectura
full-stack: backend FastAPI, autentificare JWT, integrarea API-Football
și interfața React, inclusiv fluxul complet al unui utilizator care
generează o predicție, testarea și manualul utilizatorului.

\textbf{Concluzii și direcții viitoare} sintetizează rezultatele
obținute, limitările sistemului și direcțiile de extindere, în special
tranziția către predicție în timp real.

\bigskip
\noindent\hrulefill
\bigskip

\noindent\textbf{Declarație privind utilizarea instrumentelor AI.}
Pe parcursul elaborării acestei lucrări au fost utilizate instrumente
bazate pe inteligență artificială (ChatGPT / GitHub Copilot) cu rol de
asistență în redactarea și reformularea unor pasaje de text, precum și
în generarea unor fragmente de cod auxiliar. Toate ideile, deciziile de
proiectare, implementarea sistemului și interpretarea rezultatelor îmi
aparțin în întregime. Responsabilitatea pentru conținutul lucrării este
exclusiv a autorului.